% !TEX program = xelatex
% ============================================================
%  SY02-47  射影曲线上的几何 Langlands
%  Sam Raskin 讲授，Lin Chen 笔记，2025 年 7 月 7--11 日
%  编译：xelatex 三遍
% ============================================================
\documentclass[11pt,a4paper]{article}
% ============================================================
%  SY02-47  Raskin, Geometric Langlands for Projective Curves
%  中译本 preamble
% ============================================================
\usepackage[UTF8,fontset=windows,heading=false]{ctex}
\usepackage{amsmath,amssymb,amsthm,mathtools,bm}
\usepackage{array}
\usepackage[margin=2.6cm]{geometry}
\usepackage{xcolor}
\usepackage[hidelinks]{hyperref}
\usepackage{tikz}
\usepackage{tikz-cd}
\usetikzlibrary{decorations.pathreplacing}
\numberwithin{equation}{section}

% ---------- 版式 ----------
\linespread{1.18}
\setlength{\parskip}{0.25em}
\setlength{\marginparwidth}{2.0cm}
\renewcommand{\proofname}{证明}
\renewcommand{\qedsymbol}{$\square$}

% ---------- 编号环境：同一节共用一枚计数器 ----------
\theoremstyle{definition}
\newtheorem{theorem}{定理}[subsection]
\newtheorem{proposition}[theorem]{命题}
\newtheorem{lemma}[theorem]{引理}
\newtheorem{corollary}[theorem]{推论}
\newtheorem{conjecture}[theorem]{猜想}
\newtheorem{definition}[theorem]{定义}
\newtheorem{example}[theorem]{例}
\newtheorem{remark}[theorem]{注}
\newtheorem{goal}[theorem]{目标}
\newtheorem{construction}[theorem]{构造}
\newtheorem{exercise}[theorem]{练习}
\newtheorem{warning}[theorem]{警告}
\newtheorem{question}[theorem]{问题}
\newtheorem{fact}[theorem]{事实}

% ---------- 原文页码锚 ----------
\newcommand{\srcpage}[1]{\marginpar{\raggedright\footnotesize\ttfamily\color{gray}p.#1}}

% ---------- 译注：\yizhu{内容}{类别与依据} ----------
\newcommand{\yizhu}[2]{\penalty10000\footnote{\textsf{【译注·#2】}#1}}

% ---------- 常用对象 ----------
\newcommand{\Bun}{\operatorname{Bun}}
\newcommand{\Jac}{\operatorname{Jac}}
\newcommand{\Sym}{\operatorname{Sym}}
\newcommand{\AJ}{\operatorname{AJ}}
\newcommand{\Rep}{\operatorname{Rep}}
\newcommand{\Shv}{\operatorname{Shv}}
\newcommand{\Vect}{\mathrm{Vect}}
\newcommand{\QCoh}{\operatorname{QCoh}}
\newcommand{\IndCoh}{\operatorname{IndCoh}}
\newcommand{\Ind}{\operatorname{Ind}}
\newcommand{\Coh}{\operatorname{Coh}}
\newcommand{\Perf}{\operatorname{Perf}}
\newcommand{\DMod}{D\text{-}\mathrm{mod}}
\newcommand{\LS}{\operatorname{LS}}
\newcommand{\Loc}{\operatorname{Loc}}
\newcommand{\CT}{\operatorname{CT}}
\newcommand{\Four}{\operatorname{Four}}
\newcommand{\coeff}{\operatorname{coeff}}
\newcommand{\Add}{\operatorname{add}}
\newcommand{\ev}{\operatorname{ev}}
\newcommand{\pr}{\mathrm{pr}}
\newcommand{\id}{\mathrm{id}}
\newcommand{\tr}{\operatorname{tr}}
\newcommand{\Aut}{\operatorname{Aut}}
\newcommand{\End}{\operatorname{End}}
\newcommand{\Hom}{\operatorname{Hom}}
\newcommand{\Map}{\operatorname{Map}}
\newcommand{\Spec}{\operatorname{Spec}}
\newcommand{\Pic}{\operatorname{Pic}}
\newcommand{\Symb}{\operatorname{Sym}}
\newcommand{\std}{\mathrm{std}}
\newcommand{\cusp}{\mathrm{cusp}}
\newcommand{\Nilp}{\mathrm{Nilp}}
\newcommand{\Ran}{\mathrm{Ran}}
\newcommand{\coarse}{\mathrm{coarse}}
\newcommand{\irred}{\mathrm{irred}}
\newcommand{\lisse}{\mathrm{lisse}}
\newcommand{\Lisse}{\operatorname{Lisse}}
\newcommand{\dR}{\mathrm{dR}}
\newcommand{\Betti}{\mathrm{Betti}}
\newcommand{\Gm}{\mathbb{G}_m}
\newcommand{\Ga}{\mathbb{G}_a}
\newcommand{\A}{\mathbb{A}}
\newcommand{\C}{\mathbb{C}}
\newcommand{\Q}{\mathbb{Q}}
\newcommand{\Z}{\mathbb{Z}}
\newcommand{\R}{\mathbb{R}}
\newcommand{\F}{\mathbb{F}}
\newcommand{\Ql}{\overline{\mathbb{Q}}_{\ell}}
\newcommand{\Fp}{\mathbb{F}_p}
\newcommand{\e}{\mathit{e}}
\newcommand{\Gd}{\check{G}}
\newcommand{\Td}{\check{T}}
\newcommand{\Bd}{\check{B}}
\newcommand{\Pd}{\check{P}}
\newcommand{\Md}{\check{M}}
\newcommand{\Nd}{\check{N}}
\newcommand{\LG}{\mathrm{L}G}
\newcommand{\LpG}{\mathrm{L}^{+}G}
\newcommand{\LB}{\mathrm{L}B}
\newcommand{\LT}{\mathrm{L}T}
\newcommand{\LU}{\mathrm{L}U}
\newcommand{\cF}{\mathcal{F}}
\newcommand{\cO}{\mathcal{O}}
\newcommand{\cL}{\mathcal{L}}
\newcommand{\cE}{\mathcal{E}}
\newcommand{\cD}{\mathcal{D}}
\newcommand{\cN}{\mathcal{N}}
\newcommand{\cP}{\mathcal{P}}
\newcommand{\cT}{\mathcal{T}}
\newcommand{\cZ}{\mathcal{Z}}
\newcommand{\gc}{\mathfrak{g}}
\newcommand{\bc}{\mathfrak{b}}
\newcommand{\pc}{\mathfrak{p}}
\newcommand{\mc}{\mathfrak{m}}
\newcommand{\nc}{\mathfrak{n}}
\newcommand{\tc}{\mathfrak{t}}
\newcommand{\dd}{\mathrm{d}}
\newcommand{\ee}{\mathrm{e}}
\newcommand{\one}{\mathbf{1}}


\title{射影曲线上的几何 Langlands\\[0.4em]
\large 讲义与译注}
\author{Sam Raskin 讲授\\[0.2em]
Lin Chen 笔记\\[0.6em]
\small 中文译注本}
\date{2025 年 7 月 7--11 日}

\begin{document}
\maketitle

\noindent\small
\textbf{底本}：Sam Raskin (lecturer), Lin Chen (notes), \emph{Notes for ``Geometric Langlands for Projective Curves''}, July 07--11, 2025, 18 pp.
原 PDF 由 pdfTeX 于 2025-07-20 排版；公开副本见
\url{https://www.math.columbia.edu/~okounkov/Sam_lecture_notes.pdf}。
PDF 页与印刷页一致，译文边注中的 \texttt{p.N} 即指此页。

\medskip
\noindent\small
\textbf{译注约定}：原文含糊、省略或疑似笔误之处，以脚注标出。脚注以「【译注·类别】」起首，类别只取
\emph{术语背景}、\emph{原文省略}、\emph{原文笔误}、\emph{记号}、\emph{背景直觉}、\emph{含糊点拨}
六种，并注明依据等级：
A 为原始论文、经典教材或作者本人文献；B 为 nLab、Stacks Project、ProofWiki、EMS、MathWorld、DLMF
等可追溯的百科或正式讲义；I 为本库既有译注与术语表；无外部依据者写「译者按（上下文推断）」。
译文不补写原文所无的论证；原文的笔误照录，另加批注给出正确形式。

\medskip
\noindent\small
\textbf{记号}：原文用同一个 \(e\) 兼指系数域与单位对象。译文以 \(\e\) 记系数域，以 \(\one\) 记单位对象
（平凡的秩 \(1\) 局部系统或张量单位），以免混淆。对偶群统一写 \(\Gd\)。\(\mathrm{Bun}_G\) 等模叠记号按原文。

\normalsize
\newpage
\setcounter{tocdepth}{2}
{\small
\tableofcontents}
\newpage

% ============================================================
%  第一讲
% ============================================================
\section{第一讲}

\subsection{层论}
\srcpage{2}

本课程始终固定以下设定：

\begin{itemize}
\item 一个代数闭域 \(k=\bar k\)，称为\emph{几何域}；
\item 一个系数域为 \(\e\) 的层论，满足 \(\operatorname{char}(\e)=0\)。
\end{itemize}

下面是层论的几个例子。

\begin{example}[Betti 情形]\label{ex:betti}
取 \(k=\C\)。令 \(\Shv(Y)=\Shv(Y,\e)\) 为 \(Y(\C)\) 上（复拓扑下）\(\e\)-向量空间层的范畴，
并令 \(\Lisse(Y)\subset\Shv(Y)\) 为有限秩局部常值层构成的全子范畴。
\end{example}

\begin{example}[\(\ell\)-进/étale 情形]\label{ex:etale}
取素数 \(\ell\neq\operatorname{char}(k)\)，并令 \(\e=\Ql\)。令 \(\Shv(Y)\) 为 \(Y\) 上 \(\ell\)-进层的范畴，
并令 \(\Lisse(Y)\subset\Shv(Y)\) 为 lisse 层构成的全子范畴。
\end{example}

\begin{example}[de Rham 情形]\label{ex:dr}
取 \(\operatorname{char}(k)=0\)，并令 \(\e=k\)。令 \(\Shv(Y)=\DMod(Y)\) 为 \(Y\) 上 \(D\)-模的范畴，
并令 \(\Lisse(Y)\subset\Shv(Y)\) 为带平坦联络的有限秩向量丛构成的全子范畴。
\end{example}

\begin{remark}\label{rem:rh}
当 \(k=\C\) 且 \(Y\) 为光滑簇时，有等价
\[
\Lisse_{\dR}(Y)\xrightarrow{\ \simeq\ }\Lisse_{\Betti}(Y,\C),
\]
它把 \((E,\nabla)\) 送到 de Rham 复形
\([E\to E\otimes\Omega^1_Y\to E\otimes\Omega^2_Y\to\cdots]\)。
\yizhu{``lisse'' 是法语「光滑」的借词，在层论中与「局部常值」同义；本译本保留原文，不另造中文。
Betti、de Rham、étale 三种情形的关系，通常称为 Riemann--Hilbert 对应。}{术语背景·B：nLab ``local system''（Last revised 2025-04-10）；A：Deligne, \emph{Équations différentielles à points singuliers réguliers}, LNM 163 (1970)}
\end{remark}

\subsection{几何类域论，亦称「交换几何 Langlands」}

设 \(X\) 为 \(k\) 上亏格为 \(g\) 的光滑射影曲线。考虑 \(X\) 上线丛的模叠 \(\Bun_{\Gm}\)。
固定基点 \(x_0\in X\)，有识别
\[
\Bun_{\Gm}\simeq\Jac(X)\times\Z\times B\Gm.
\]
其中：

\begin{itemize}
\item \(\Jac(X)\) 是 \(X\) 的雅可比簇，它参数化 \(X\) 上次数为 \(0\)、且在 \(x_0\) 处带平凡化的线丛；
\item 因子 \(\Z\) 记录线丛的次数；
\item 因子 \(B\Gm\) 记录线丛的自同构，自同构群为 \(k^\times\)。
\end{itemize}

\begin{goal}\label{goal:gl1}
我们要作如下构造：

\begin{itemize}
\item （输入）\(X\) 上一个秩 \(1\) 局部系统（即 lisse 层）\(\sigma\)；
\item （输出）\(\Bun_{\Gm}\) 上一个 lisse 层 \(\chi_\sigma\)，并附带以下数据：
\begin{itemize}
\item （Hecke 性质）对任意 \(L\in\Bun_{\Gm}\)、\(x\in X\)，记 \(L(x):=L\otimes_{\cO}\cO(x)\)，有一个同构
\[
\chi_\sigma|_{L(x)}\simeq\chi_\sigma|_L\otimes\sigma_x,
\]
它代数地依赖于 \((L,x)\)；
\item （归一化）一个同构
\[
\chi_\sigma|_{\cO_X}\simeq\one.
\]
\end{itemize}
\end{itemize}
\end{goal}

下面给出 \(\chi_\sigma\) 的两种构造。

\begin{construction}\label{con:gl1-betti}
为简单起见，先在 Betti 情形工作。此时知道 \(\sigma\) 等价于知道一个同态
\[
\pi_1(X,x_0)\to\e^\times.
\]
把基本群换成它的交换化，得到
\[
\pi_1(X,x_0)^{\mathrm{ab}}\simeq H_1(X,\Z)\simeq H^1(X,\Z),
\]
最后一个同构来自 Poincaré 对偶。于是 \(\sigma\) 给出一个同态
\begin{equation}\label{eq:1.1}
H_1(X,\Z)\to\e^\times.
\end{equation}
另一方面，短正合列
\[
0\to\Z\to\cO\xrightarrow{\ \exp\ }\cO^\times\to 0
\]
诱导出同构
\[
\Jac(X)\simeq H^1(X,\cO)/H^1(X,\Z),
\]
从而
\[
\pi_1(\Jac(X))\simeq H_1(X;\Z).
\]
结合 \eqref{eq:1.1}，得到同态
\[
\pi_1(\Jac(X))\to\e^\times,
\]
它给出 \(\Jac(X)\) 上一个秩 \(1\) 局部系统 \(\chi_\sigma|_{\Jac(X)}\)。现在定义 \(\chi_\sigma\)，使其在
\(\Jac(X)\times\{d\}\times B\Gm\) 上的限制为
\[
\chi_\sigma|_{\Jac(X)}\boxtimes\sigma_{x_0}^{\otimes d}\boxtimes\one_{B\Gm}.
\]
\end{construction}

\begin{exercise}\label{ex:gl1-hecke}
验证 \(\chi_\sigma\) 满足 Hecke 性质。
\end{exercise}

\begin{construction}[Deligne]\label{con:gl1-deligne}
这个构造对任意层论都适用。

对线丛 \(L\in\Bun_{\Gm}\)，选取一个有理截面就给出同构 \(\cO(D)\simeq L\)，其中 \(D=\sum n_i x_i\) 是 \(X\) 上的除子。
我们对 \(\chi_\sigma\) 的公理要求
\[
\chi_\sigma|_L\simeq\bigotimes_i\sigma_{x_i}^{\otimes n_i}.
\]
由此可见 \(\chi_\sigma\) 是「超定的」；需要回答的问题是：这个向量空间为什么与 \(D\) 的选取无关？

为处理这个问题，考虑 \(X\) 的 \(d\) 次对称幂 \(\Sym^d(X)\)，它是 \(X\) 上次数为 \(d\) 的有效除子的模空间。
等价地，它分类 \((L,s)\)，其中 \(L\) 是次数为 \(d\) 的线丛，\(s\in\Gamma(X,L)\) 是非零截面。有映射
\[
\AJ_d\colon\Sym^d(X)\to\Bun^d_{\Gm},
\]
它忘掉截面 \(s\)。这就是 Abel--Jacobi 映射。该映射在 \(L\) 处的纤维是 \(\Gamma(X,L)\setminus 0\)。
\yizhu{原文说 \(\AJ_d\) 在 \(L\) 处的纤维是 \(\Gamma(X,L)\setminus 0\)。严格说，纤维是 \(\Gamma(X,L)\) 的射影化
\(\mathbb P\Gamma(X,L)\)：\(\AJ_d\) 按除子定义，两个非零截面对应同一个除子，当且仅当它们相差一个非零常数倍。
维数上也能看出：\(\dim\Sym^d(X)=d\)，而 \(\dim\Gamma(X,L)=d+1-g\)，若纤维是 \(\Gamma\setminus0\)，总维数会是
\(2d+1-g\)，与 \(d\) 矛盾；\(\dim\mathbb P\Gamma(X,L)=d-g\) 才相符。注意 \(\Gamma(X,L)\setminus 0\) 与
\(\mathbb P\Gamma(X,L)\) 在 \(d>2g-2\) 时都单连通（前者同伦等价于
\(\mathbb C^{d+1-g}\setminus0\simeq S^{2(d+1-g)-1}\)），单连通性不能区分二者；区别在于维数与「纤维记录除子而非截面」。}{原文笔误·A：纤维按除子定义、维数 \(d-g\)；B：Feng, \emph{Geometric class field theory} 讲义 §2}

当 \(d>2g-2\) 时，由 Riemann--Roch，
\[
\dim\Gamma(X,L)=d+1-g,
\]
这说明 \(\AJ_d\) 是光滑的。注意 \(\AJ_d\) 的纤维单连通。

现在对给定的 \(\sigma\)，可以作出它的 \(d\) 次对称幂，即 \(\Sym^d(X)\) 上一个秩 \(1\) 局部系统，公式为
\begin{equation}\label{eq:1.2}
\sigma^{(d)}:=\Add_{d,*}\bigl(\sigma\boxtimes\cdots\boxtimes\sigma\bigr)^{S_d}.
\end{equation}
这里 \(\Add_d\colon X^d\to\Sym^d(X)\)，\(S_d\) 是对称群。我们的公理要求
\[
\chi_\sigma|_{\Sym^d(X)}\simeq\sigma^{(d)}.
\]
由前述单连通性，当 \(d\gg 0\) 时，\(\sigma^{(d)}\) 必须下降为 \(\Bun^d_{\Gm}\) 上的 lisse 层。
于是可以先对 \(d\gg 0\) 定义 \(\chi_\sigma|_{\Bun^d_{\Gm}}\)，再利用 Hecke 性质延拓到整个 \(\Bun_{\Gm}\)。
\end{construction}

\begin{remark}\label{rem:sym}
在 \eqref{eq:1.2} 中，\(\Add_{d,*}(\sigma\boxtimes\cdots\boxtimes\sigma)\) 在 \(D\) 处的纤维为
\[
\bigoplus_{D=\sum x_i}\bigotimes_i\sigma_{x_i},
\]
其中直和按把 \(D\) 写成 \(\sum x_i\) 的所有方式取。现在取 \(S_d\)-不变量，就消去了这种重复，
使纤维变成 \(\bigotimes_i\sigma_{x_i}\)。
\yizhu{这里 \(D=\sum x_i\) 应理解为把重数拆开的写法：若 \(D=\sum_i n_i x_i\)，则不变纤维是
\(\bigotimes_i\sigma_{x_i}^{\otimes n_i}\)。原文简写为 \(\bigotimes_i\sigma_{x_i}\)，未标重数。}{记号·I：本工程 §1.2 统一约定}
\end{remark}

\subsection{非交换理论}
\srcpage{4}

现在取 \(G=\mathrm{PGL}_2\)。目标变为：

\begin{itemize}
\item （输入）\(X\) 上一个 \(\mathrm{SL}_2\)-局部系统 \(\sigma\)。我们把 \(\sigma\) 看作一个秩 \(2\) 局部系统，
并带有一个 \(\wedge^2\sigma\) 的平凡化；
\item （输出）\(\Bun_{\mathrm{PGL}_2}\) 上一个满足 Hecke 性质的层 \(\cF_\sigma\)（Hecke 性质将在后面的讲次解释）。
\end{itemize}

不过，我们不再要求 \(\cF_\sigma\) 是 lisse 的；只要求它是 \emph{preverse} 的。理由将在后面的讲次说明。
\yizhu{原文如此，应为 \emph{perverse}，即「反常的」。这是原文的拼写笔误；译文照录原文，正确形式见此注。}{原文笔误·A：后文 §3.1 的 ``perverse but not lisse'' 可对照}

上文中，\(\Bun_{\mathrm{PGL}_2}\) 是 \(X\) 上 \(\mathrm{PGL}_2\)-丛的模叠。按定义，\(X\) 上的 \(\mathrm{PGL}_2\)-丛
是秩 \(2\) 向量丛模去「张量一个线丛」的歧义。换言之，\(\cE\) 与 \(\cE\otimes\cL\) 给出 \(\Bun_{\mathrm{PGL}_2}\) 中同一个点。
注意
\[
\deg(\cE\otimes\cL)=\deg(\cE)+2\deg(\cL).
\]
因此有良定义的映射
\[
\Bun_{\mathrm{PGL}_2}\xrightarrow{\ \deg\ }\Z/2,
\]
以及分解
\[
\Bun_{\mathrm{PGL}_2}\simeq\Bun_{\mathrm{PGL}_2}^{\mathrm{even}}\sqcup\Bun_{\mathrm{PGL}_2}^{\mathrm{odd}}.
\]

\begin{figure}[htbp]
\centering
\begin{tikzpicture}[scale=1.05, every node/.style={font=\small}]
  % even tower
  \fill[gray!12] (0,0) rectangle (3.4,1.3);
  \draw[thick] (0,0) rectangle (3.4,5.4);
  \foreach \y in {1.3,2.7,4.1} \draw[dashed] (0,\y)--(3.4,\y);
  \node at (1.7,0.65) {\(\Bun^{\mathrm{ss}}_{\mathrm{PGL}_2}\)};
  \node at (1.7,2.0) {\(\Bun^2_B\)};
  \node at (1.7,3.4) {\(\Bun^4_B\)};
  \node at (1.7,4.75) {\(\vdots\)};
  \node at (1.7,-0.55) {\(\Bun^{\mathrm{even}}_{\mathrm{PGL}_2}\)};
  % odd tower
  \begin{scope}[xshift=6.6cm]
  \fill[gray!12] (0,0) rectangle (3.4,1.3);
  \draw[thick] (0,0) rectangle (3.4,5.4);
  \foreach \y in {1.3,2.7,4.1} \draw[dashed] (0,\y)--(3.4,\y);
  \node at (1.7,0.65) {\(\Bun^{\mathrm{ss}}_{\mathrm{PGL}_2}\)};
  \node at (1.7,2.0) {\(\Bun^1_B\)};
  \node at (1.7,3.4) {\(\Bun^3_B\)};
  \node at (1.7,4.75) {\(\vdots\)};
  \node at (1.7,-0.55) {\(\Bun^{\mathrm{odd}}_{\mathrm{PGL}_2}\)};
  \end{scope}
  % 标注：底部阴影为开子叠 Bun^{ss}；整个有限并 U_n 也是开子叠
  \draw[decorate,decoration={brace,amplitude=4pt,mirror}] (-0.55,0)--(-0.55,1.3)
    node[midway,left=4pt,rotate=90,font=\footnotesize]{开};
  \draw[decorate,decoration={brace,amplitude=4pt}] (10.55,0)--(10.55,5.4)
    node[midway,right=4pt,rotate=90,font=\footnotesize]{\(U_n\)};
\end{tikzpicture}
\caption{\(\Bun_{\mathrm{PGL}_2}\) 的示意图（据原文手绘图重绘）：两个连通分支
\(\Bun^{\mathrm{even}}\) 与 \(\Bun^{\mathrm{odd}}\)，阴影部分为半稳定轨迹
\(\Bun^{\mathrm{ss}}\)（开子叠），上方依次为 \(\Bun^d_B\) 分层；取到第 \(n\) 层时，
\(\Bun^{\mathrm{ss}}\) 与这些分层的并记为 \(U_n\)，它也是开子叠。}
\label{fig:bunpgl2}
\end{figure}

上图是 \(\Bun_{\mathrm{PGL}_2}\) 的示意图。需要作几点说明：

\begin{itemize}
\item \(\Bun_{\mathrm{PGL}_2}^{\mathrm{even}}\) 与 \(\Bun_{\mathrm{PGL}_2}^{\mathrm{odd}}\) 都连通。
\item 半稳定轨迹由满足下述条件的 \(\cE\) 组成：对任意线子丛 \(L\subset\cE\)，都有
\[
\deg(L)\le\deg(\cE)/2.
\]
注意这个条件在张量线丛下不变。它是开条件，因此定义出开子叠
\[
\Bun_{\mathrm{PGL}_2}^{\mathrm{ss}}\subset\Bun_{\mathrm{PGL}_2}.
\]
\item 设 \(B\subset\mathrm{PGL}_2\) 为标准 Borel 子群，\(\Bun_B\) 为 \(X\) 上 \(B\)-丛的模叠。按定义，
它分类短正合列
\[
[0\to L\to\cE\to\cO\to 0].
\]
有映射
\[
\Bun_B\xrightarrow{\ \deg\ }\Z,
\]
把上述正合列送到 \(\deg(L)\)。这给出分解
\[
\Bun_B\simeq\bigsqcup_{d\in\Z}\Bun^d_B.
\]
当 \(d\ge 1\) 时，映射
\[
\Bun^d_B\to\Bun_{\mathrm{PGL}_2}
\]
是局部闭嵌入。
\item 对 \(n\ge 0\)，有开子叠 \(U_n\subset\Bun_{\mathrm{PGL}_2}\)，它是 \(\Bun^{\mathrm{ss}}_{\mathrm{PGL}_2}\)
与 \(\Bun^d_B\)（\(1\le d\le n\)）的像的不交并。特别地，\(\Bun_{\mathrm{PGL}_2}\) 不是拟紧的。
\yizhu{\(\Bun_{\mathrm{PGL}_2}\) 不是拟紧的：它含有无穷多个互不相同的分层 \(\Bun^d_B\)（\(d\ge 1\)），
而 \(U_n\) 只包含前 \(n\) 个。这一点在 §2.1 的尖点延拓中会用到：尖点层虽然可以「干净地」从某个 \(U_n\)
延拓，但 \(n\) 的选取依赖 \(g\)，不依赖层本身。}{原文省略·A：Drinfeld--Gaitsgory, \emph{Geometric constant term functor(s)}, Selecta Math. 22 (2016), §2；I：本工程 §2.1}
\end{itemize}

\begin{exercise}\label{ex:connected}
证明 \(\Bun_{\mathrm{PGL}_2}^{\mathrm{even}}\) 与 \(\Bun_{\mathrm{PGL}_2}^{\mathrm{odd}}\) 都连通。
\end{exercise}

\begin{exercise}\label{ex:inj}
证明：对 \(d\ge 0\)，映射 \(\Bun^d_B\to\Bun_{\mathrm{PGL}_2}\) 在 \(k\)-点上为单射。
\end{exercise}

\begin{exercise}\label{ex:dim}
求 \(\Bun^d_B\) 与 \(\Bun_{\mathrm{PGL}_2}\) 的维数，并推出存在奇的半稳定 \(\mathrm{PGL}_2\)-丛。
\end{exercise}

\begin{exercise}\label{ex:smooth}
对 \(d\ll 0\)，证明 \(\Bun^d_B\to\Bun_{\mathrm{PGL}_2}\) 是光滑的。
\end{exercise}

% ============================================================
%  第二讲
% ============================================================
\section{第二讲}

上次我们说明：给定 \(X\) 上一个不可约 \(\mathrm{SL}_2\)-局部系统，我们希望构造 \(\Bun_{\mathrm{PGL}_2}\) 上一个对应的层 \(\cF_\sigma\)，
并列出它应当满足的一些性质（当时尚未精确化）。这些性质包括：

\begin{itemize}
\item \(\cF_\sigma\) 应当是反常的；
\item \(\cF_\sigma|_{\Bun^{\mathrm{even}/\mathrm{odd}}_{\mathrm{PGL}_2}}\) 作为反常层应当不可约；
\item \(\cF_\sigma\) 应当是尖点的；
\item \(\cdots\)
\end{itemize}

在上述诸性质中，尖点性是「\(\sigma\) 不可约」这一条件的对应物；这是 Langlands 哲学的一部分。

\subsection{尖点性}
\srcpage{6}

\begin{definition}\label{def:ct}
考虑对应
\[
\begin{tikzcd}
& \Bun_B \arrow[dl,"p"'] \arrow[dr,"q"] & \\
\Bun_{\mathrm{PGL}_2} & & \Bun_{\Gm}
\end{tikzcd}
\]
其中 \(p\) 把正合列 \([0\to L\to\cE\to\cO\to 0]\in\Bun_B\) 送到 \(\cE\in\Bun_{\mathrm{PGL}_2}\)，
\(q\) 把它送到 \(L\in\Bun_{\Gm}\)。定义\emph{常数项函子}
\[
\CT_*:=q_*p^!\colon\Shv(\Bun_{\mathrm{PGL}_2})\to\Shv(\Bun_{\Gm}).
\]
对每个 \(d\in\Z\)，把 \(\Bun_B\) 换成 \(\Bun^d_B\)，类似地得到函子
\[
\CT^d_*\colon\Shv(\Bun_{\mathrm{PGL}_2})\to\Shv(\Bun^d_{\Gm}).
\]
\end{definition}

\begin{definition}\label{def:cuspidal}
称 \(F\in\Shv(\Bun_G)\) 是\emph{尖点的}，若 \(\CT_*(F)\simeq 0\)。
\yizhu{「尖点」这个译名来自自守表示：cuspidal automorphic representation 的中文惯译为「尖点自守表示」。
在几何情形，常数项函子 \(\CT_*\) 是模形式常数项 \(a_0\) 的层论类比；「尖点」即常数项为零。
本文把 ``cuspidal'' 统一译为「尖点」（作形容词时写「尖点层」），不译作「尖端的」。}{术语背景·I：本库既有译名；A：Frenkel, \emph{Langlands Correspondence for Loop Groups}, §3}
\end{definition}

\begin{remark}\label{rem:cusp-all-d}
\(F\) 是尖点的，当且仅当对任意 \(d\in\Z\)，都有 \(\CT^d_*(F)\simeq 0\)。
\end{remark}

\begin{remark}\label{rem:cusp-general}
对一般约化群，我们用所有真抛物子群及其 Levi 商来定义尖点性。
\end{remark}

回忆上一讲的图景 \eqref{fig:bunpgl2}（原文误引作 (1.3)）。下述结果表明，尖点性蕴含「在无穷远处」的消没。

\begin{proposition}\label{prop:cusp-vanish}
若 \(F\) 是尖点的，则它对 \(\Bun^d_B\) 的 \(!\)-限制在 \(d>2g-2\) 时为零。
\end{proposition}

\begin{proof}[证明梗概]
考虑映射 \(q_d\colon\Bun^d_B\to\Bun^d_{\Gm}\)。若 \(d>2g-2\)，则 \(\Bun^d_B\) 中任一点对应的正合列
\(0\to L\to B\to\cO\to 0\) 都分裂\yizhu{原文此处写 \(0\to L\to B\to\cO\to0\)，其中 \(B\) 是 \(\cE\) 之误；同页下文与脚注均用 \(\cE\)。
译文照录原文，正确形式为 \(\cE\)。}{原文笔误·A：同页自相矛盾}。于是 \(q_d\) 在 \(L\in\Bun^d_{\Gm}\) 处的纤维可以等同于
分类叠 \(B H^1(X,L)\)。

回忆：由于 \(\Ga\) 可缩，有等价
\[
\pi^*\colon\Vect\rightleftarrows\Shv(B\Ga):\pi_*,
\]
其中 \(\pi\colon B\Ga\to\mathrm{pt}\) 是投影。类似地，由于 \(H^1(X,L)\) 可缩，函子
\[
q_{d,*}\colon\Shv(\Bun^d_B)\to\Shv(\Bun^d_{\Gm})
\]
在 \(d>2g-2\) 时是等价。于是 \(\CT^d_*(F)\simeq 0\) 蕴含 \(p^!_d(F)\simeq 0\)，这正是所要的。
\yizhu{原文正文说 \(q_d\) 的纤维是 \(B H^1(X,L)\)，脚注则说纤维可典范地等同于向量叠 \(R\Gamma(X,L)[1]\)。
脚注才是正确的导出陈述。当 \(d>2g-2\) 时，\(H^1(X,L)=0\) 而 \(H^0(X,L)\neq 0\)，故
\(R\Gamma(X,L)[1]\simeq H^0(X,L)[1]\)，纤维是 \(B H^0(X,L)\)，不是 \(B H^1(X,L)\)。原文紧接着的
「因为 \(H^1(X,L)\) 可缩」也应改为「因为 \(H^0(X,L)\) 是向量空间」：向量群的有限维表示平凡，故
\(\Shv(B H^0)\simeq\Vect\)，\(q_{d,*}\) 仍是等价。}{原文笔误·A：脚注与正文不一致；Drinfeld--Gaitsgory, \emph{Geometric constant term functor(s)}, Selecta Math. 22 (2016), §2}
\end{proof}

我们还可以定义另一版本的常数项函子。

\begin{definition}\label{def:ct!}
定义 \(\CT^!\) 为函子
\[
q_!p^*\colon\Shv(\Bun_G)\to\Shv(\Bun_T).
\]
\end{definition}

\begin{remark}\label{rem:ct-duality}
形式地说，若 \(F\) 可构造，则 \(\CT_*(F)\simeq D\CT^!D(F)\)，其中 \(D\) 是 Verdier 对偶。
\end{remark}

\begin{theorem}[Drinfeld--Gaitsgory，「第二伴随」]\label{thm:second-adjoint}
有典范等价
\[
\CT^d_*\simeq\mathrm{inv}\circ\CT^{-d}_!,
\]
其中 \(\mathrm{inv}\) 是由同构 \(\Bun^d_{\Gm}\simeq\Bun^{-d}_{\Gm}\)、\(L\mapsto L^{-1}\) 诱导的函子。
\end{theorem}

\begin{corollary}\label{cor:dual-cusp}
若 \(F\) 可构造且尖点，则 \(DF\) 也尖点。
\end{corollary}

\begin{corollary}\label{cor:star-vanish}
若 \(F\) 是尖点的，则它对 \(\Bun^d_B\) 的 \(*\)-限制在 \(d>2g-2\) 时为零。
\end{corollary}

\begin{proof}
由 Drinfeld--Gaitsgory 定理，\(\CT^{-d}_*(F)\simeq 0\) 蕴含 \(\CT^d_!(F)\simeq 0\)。
与命题 \ref{prop:cusp-vanish} 的证明类似，后者蕴含 \(p^*_d(F)\simeq 0\)（\(d>2g-2\)）。
\end{proof}

\begin{corollary}\label{cor:extend}
存在开子叠 \(U_n\subset\Bun_G\)，它是有限个分层的并，使得任意尖点层 \(F\) 都能从 \(U_n\)
同时作 \(!\)-延拓与 \(*\)-延拓。
\yizhu{原文脚注 2 只说 \(F\) 从 \(U_n\)「干净地延拓」（cleanly extended）。译者按（上下文推断，无外部依据）：
「干净」应指 \(!\)-延拓与 \(*\)-延拓重合，因而延拓既不丢失也不增加信息；原文未给定义。}{原文省略·A：原文脚注 2}
\end{corollary}

\begin{remark}\label{rem:n}
事实上可以取 \(n=\max\{0,2g-2\}\)。
\end{remark}

\begin{exercise}\label{ex:second-adjoint-constant}
对 \(\Bun_G\) 上的常层验证定理 \ref{thm:second-adjoint}。
\end{exercise}

\begin{exercise}\label{ex:opposite-borel}
设 \(B^-\subset G\) 为标准反 Borel 子群，并用 \(B^-\) 代替 \(B\) 定义 \(\CT^-_*\)、\(\CT^-_!\)。
证明：定理 \ref{thm:second-adjoint} 等价于 \(\CT^-_!\simeq\CT_*\)。挑战：构造自然变换
\(\CT^-_!\to\CT_*\) 与 \(\CT_*\to\CT^-_!\)。
\end{exercise}

\begin{exercise}\label{ex:cusp-exist}
（挑战）对 \(g\ge 2\)，证明尖点层存在。
\end{exercise}

\subsection{系数}
\srcpage{7}

为了陈述对 \(\cF_\sigma\) 的其他预期，我们借用模形式理论（亦即 \(\mathrm{PGL}_2\) 的自守形式理论）中的一些动机。

回忆：权为 \(k\)、级为 \(1\) 的 holonomic 模形式
\yizhu{权 \(k\) 由 \(\sigma\) 在阿基米德位处的行为决定。}{原文省略·A：原文脚注 3}
是一个和
\[
f=\sum_{n\ge 0}a_n q^n,
\]
它在 \(|q|<1\) 时收敛，并且对任意 \(g=\begin{pmatrix}a&b\\c&d\end{pmatrix}\in\mathrm{SL}_2(\Z)\)，满足
\[
f(g\cdot\tau)=(c\tau+d)^k f(\tau),
\]
其中 \(q=\exp(2\pi i\tau)\)，\(\operatorname{Im}(\tau)>0\)。
\yizhu{原文把模形式写成 ``(holonomic) modular form''，译文照录。按上下文应为 ``holomorphic''（全纯），
``holonomic'' 是笔误（holonomic 是 \(D\)-模的术语，与模形式的全纯性无关）。}{原文笔误·A：全文术语一致性}

Langlands 猜想断言：对任意（奇、不可约、非分歧的）Galois 表示
\[
\sigma\colon\text{``}\mathrm{Gal}_{\Q}\text{''}\to\mathrm{SL}_2(\text{``}\C\text{''}),
\]
存在模形式 \(f_\sigma=\sum a_n q^n\)，使得

\begin{itemize}
\item （尖点性）常数项 \(a_0\) 为零；
\item \(a_p=\operatorname{tr}(\sigma(\mathrm{Fr}_p))\)，其中 \(\mathrm{Fr}_p\in\mathrm{Gal}_{\Q}\) 是素数 \(p\) 的 Frobenius 共轭类；
\item \(a_1=1\)，当 \((n,m)=1\) 时 \(a_{nm}=a_n a_m\)，并且
\begin{equation}\label{eq:2.1}
a_{p^{n+1}}=a_p a_{p^n}-p^{k-1}a_{p^{n-1}}.
\end{equation}
\end{itemize}

注意上述条件唯一确定 \(f_\sigma\)。

我们的目标是为 \(\Bun_{\mathrm{PGL}_2}\) 上的层提供 \(a_n\) 的类比，其中 \(a_n\) 是模形式的第 \(n\) 个 Fourier 系数。

\begin{itemize}
\item 我们把 \(n\)（\(\operatorname{Spec}\Z\) 上的有效除子）换成 \(X\) 上的有效除子 \(D\)；
\item 按照 Grothendieck 关于层--函数对应的哲学，我们把数 \(a_n(f)\) 换成向量空间 \(\coeff_D(F)\in\Vect\)。
\end{itemize}

先考虑 \(D=0\) 的情形，它类比于 \(a_1(f)\)。回忆
\[
a_n(f)=\int_{\R/\Z}f(\tau)\exp(-2\pi i n\tau)\,d\tau.
\]
这里我们把 \(\R\) 认同为 \(\mathrm{SL}_2(\R)\) 的标准 Borel 子群的单幂根。这促使我们考虑 \(\Bun_N\)，
其中 \(N\) 是 \(\mathrm{PGL}_2\) 的标准 Borel 子群的单幂根。按定义，\(\Bun_N\) 是扩张
\([0\to\cO\to\cE\to\cO\to 0]\) 的模叠。

由于一些微妙的原因，我们需要考虑一个变体
\[
\Bun^\Omega_N:=\Bun_B\times_{\Bun_{\Gm}}\{\Omega\},
\]
其中 \(\Omega\in\Bun_{\Gm}\) 是 \(X\) 上典范线丛对应的点。注意作为向量叠，有
\[
\Bun^\Omega_N\simeq R\Gamma(X,\Omega)[1].
\]
现在考虑对应
\[
\begin{tikzcd}
& \Bun^\Omega_N \arrow[dl,"p_N"'] \arrow[dr,"\phi"] & \\
\Bun_G & & H^1(X,\Omega)\simeq\A^1
\end{tikzcd}
\]
其中同构 \(H^1(X,\Omega)\simeq\A^1\) 来自 Serre 对偶。

\begin{definition}\label{def:coeff0}
定义
\[
\coeff_0(F):=\mathrm{C}^\bullet_c\bigl(\Bun^\Omega_N,\ p_N^!(F)\overset{!}{\otimes}\phi^!(\exp)\bigr),
\]
其中 \(\exp\in\Shv(\A^1)\) 是「指数层」（下面解释）。
\yizhu{\(\mathrm{C}^\bullet_c\) 是紧支集上同调。原文记作 \(\mathrm{C}^\bullet\)，未写下标；
按 \S2.2 的配对与 \(!\)-推前，此处应理解为紧支集版本。}{记号·I：本工程统一记号}
\end{definition}

\begin{example}[de Rham 情形]\label{ex:exp-dr}
取 \(\exp\) 为 \(D\)-模 \((\cO,\nabla=d-dt)\)。
\end{example}

\begin{example}[\(\ell\)-进情形]\label{ex:exp-et}
有短正合列
\[
0\to\F_p\to\Ga\xrightarrow{\ \pi\ }\Ga\to 0,
\]
其中 \(\pi\) 是 Artin--Schreier 映射 \(t\mapsto t^p-t\)。于是 \(\pi_*(\Ql)\) 带有 \(\F_p\)-作用。
固定一个非平凡特征 \(\psi\colon\F_p\to\Ql^\times\)，可以考虑 \(\pi_*(\Ql)\) 的 \(\psi\)-分量；
可以证明它是秩 \(1\) lisse 层。我们把这个层定义为 \(\exp\)。
\end{example}

\begin{remark}[Betti 情形]\label{rem:exp-betti}
在 Betti 情形，\(\exp\) 不存在；但有一些技巧可以绕开它。
\end{remark}

% ============================================================
%  第三讲
% ============================================================
\section{第三讲}

\subsection{系数（续）}
\srcpage{8}

上次我们定义了 \(\coeff_0\colon\Shv(\Bun_G)\to\Vect\)。现在对 \(X\) 上任意有效除子 \(D\) 定义 \(\coeff_D\)。

考虑叠
\[
\Bun^{\Omega(-D)}_N:=\Bun_B\times_{\Bun_{\Gm}}\{\Omega(-D)\}.
\]
它分类扩张
\[
[0\to\Omega(-D)\to\cE\to\cO\to 0].
\]
设
\[
p_{N,D}\colon\Bun^{\Omega(-D)}_N\to\Bun_G
\]
是忘掉扩张结构、只记住 \(\cE\) 的映射。考虑复合
\[
\phi_D\colon\Bun^{\Omega(-D)}_N\to H^1(X,\Omega(-D))\to H^1(X,\Omega)\simeq\A^1,
\]
其中最后一个同构来自 Serre 对偶。

\begin{definition}\label{def:coeffD}
定义
\[
\coeff_D(F):=\mathrm{C}^\bullet_c\bigl(\Bun^{\Omega(-D)}_N,\ p_{N,D}^!(F)\overset{!}{\otimes}\phi_D^!(\exp)\bigr).
\]
\end{definition}

现在可以陈述我们的目标了。

\begin{goal}\label{goal:main}
我们要作如下构造：

\begin{itemize}
\item （输入）\(X\) 上一个不可约 \(\mathrm{SL}_2\)-局部系统 \(\sigma\)；
\item （输出）一个尖点反常层 \(\cF_\sigma\in\Shv(\Bun_{\mathrm{PGL}_2})\)，它在
\(\Bun^{\mathrm{even}/\mathrm{odd}}_{\mathrm{PGL}_2}\) 上不可约，并附带以下数据：
\begin{itemize}
\item 对每个有效除子 \(D=\sum n_i x_i\)，一个同构
\begin{equation}\label{eq:3.1}
\coeff_D(\cF_\sigma)\simeq\bigotimes_i\Sym^{n_i}(\sigma_{x_i}).
\end{equation}
\end{itemize}
\end{itemize}
\end{goal}

\begin{remark}\label{rem:recursion}
\eqref{eq:3.1} 类比于上一讲的乘法性与递推公式 \eqref{eq:2.1}，因为对 \(\mathrm{SL}_2\) 的标准表示 \(V\)，有
\[
\Sym^n(V)\otimes V\simeq\Sym^{n+1}(V)\oplus\Sym^{n-1}(V).
\]
\end{remark}

上述同构 \eqref{eq:3.1} 应当代数地依赖于 \(D\)。更精确地说，对每个 \(d\ge 0\)，可以构造函子
\[
\coeff_d\colon\Shv(\Bun_G)\to\Shv(\Sym^d(X)),
\]
使得对 \(D\in\Sym^d(X)\)，\(\coeff_D\) 可以等同于复合
\[
\Shv(\Bun_G)\xrightarrow{\ \coeff_d\ }\Shv(\Sym^d(X))\xrightarrow{\ (-)|^!_D\ }\Vect.
\]
于是 \eqref{eq:3.1} 应当升级为同构
\[
\coeff_d(\cF_\sigma)\simeq\sigma^{(d)},
\]
其中
\[
\sigma^{(d)}:=\Add_{d,*}(\sigma^{\boxtimes d})^{S_d}
\]
按秩 \(1\) 情形中的方式定义。然而，\(\sigma^{(d)}\) 是反常的，但不是 lisse 的。
\yizhu{原文说 \(\sigma^{(d)}\) ``is perverse but not lisse''，这是笔误：\(\sigma^{(d)}\) 是 \(\Sym^d(X)\)
上的局部系统，按定义就是 lisse 的；这里想说的应当是「\(\cF_\sigma\) 是反常的，但不是 lisse 的」，
与秩 \(1\) 情形中 \(\chi_\sigma\) 为 lisse 形成对照。按反常层的标准归一化，\(\Sym^d(X)\)（\(d\) 维）上的
局部系统要移位 \(d\) 次才落在反常 t-结构的心内；原文断言的是 \(\cF_\sigma\)，不涉及这个移位。}{原文笔误·A：与上一段「不再要求 lisse，只要求反常」对照；B：nLab ``perverse sheaf''}

基于模形式的动机，我们可能希望函子
\[
\prod_{d\ge 0}\coeff_d\colon\Shv(\Bun_G)_{\mathrm{cusp}}\to\prod_{d\ge 0}\Shv(\Sym^d(X))
\]
是全忠实的（从而 \(\cF_\sigma\) 由上述预期唯一确定）。然而这并不成立：右端按 \(\Z_{\ge 0}\) 分裂成若干因子，
而左端只按 \(\Z/2\) 分裂。尽管如此，我们有下述断言。

\begin{proposition}\label{prop:determined}
设 \(F\) 是 \(\Bun_{\mathrm{PGL}_2}\) 上的反常尖点层，且它在
\(\Bun^{\mathrm{even}/\mathrm{odd}}_{\mathrm{PGL}_2}\) 上的限制不可约、支撑稠密。
则对任意 \(d>N\)，\(F\) 由 \(\coeff_d(F)\) 与 \(\coeff_{d+1}(F)\) 唯一确定，其中 \(N\) 是只依赖于 \(g\) 的整数。
\end{proposition}

\begin{remark}[警告]\label{rem:only-pgl2}
上述断言是 \(\mathrm{PGL}_2\) 特有的。
\end{remark}

\subsection{Fourier 变换}
\srcpage{9}

为解释命题 \ref{prop:determined}，我们回忆 Fourier（--Deligne）变换。对有限维向量空间 \(V\) 及其对偶 \(V^\vee\)，
考虑对应
\[
\begin{tikzcd}
& V\times V^\vee \arrow[dl,"\mathrm{pr}_1"'] \arrow[dr,"\mathrm{pr}_2"] & \\
V & & V^\vee
\end{tikzcd}
\]
有等价
\[
\Four\colon\Shv(V)\xrightarrow{\ \simeq\ }\Shv(V^\vee),\qquad
F\mapsto \mathrm{pr}_{2,*}\bigl(\mathrm{pr}_1^!(F)\overset{!}{\otimes}\mathrm{ev}^!(\exp)\bigr),
\]
其中 \(\mathrm{ev}\colon V\times V^\vee\to\A^1\) 是配对映射。这可以推广到基 \(S\) 上的向量丛 \(E\) 与 \(E^\vee\)，
于是有典范等价
\[
\Four\colon\Shv(E)\xrightarrow{\ \simeq\ }\Shv(E^\vee).
\]

\begin{remark}\label{rem:four-betti}
在 \(\exp\) 不存在的 Betti 情形，我们可以对单值（monodromic）层作 Fourier 变换。
\end{remark}

我们可以用 Fourier 变换与丛的几何重写 \(\coeff_d\)。回忆映射
\[
q_{-d}\colon\Bun^{-d}_B\to\Bun^{-d}_{\Gm}
\]
在 \(L\) 处的纤维等同于 \(R\Gamma(X,L)[1]\)。对 \(d>0\)，有 \(H^0(X,L)\simeq 0\)，因此 \(q_{-d}\) 是
\(\Bun^{-d}_{\Gm}\) 上的向量丛，记作 \(E_d\)。对偶丛 \(E_d^\vee\) 在 \(L\) 处的纤维为
\[
H^1(X,L)^\vee\simeq H^0(X,L^\vee\otimes\Omega).
\]
于是 \(E_d^\vee\) 是分类 \((L,s\colon L\to\Omega)\) 的模叠，其中 \(L\in\Bun^{-d}_{\Gm}\)，\(s\colon L\to\Omega\) 是映射。

注意有识别
\[
\Sym^{d+(2g-2)}(X)\to E_d^\vee\setminus\{0\},\qquad D\mapsto[\Omega(-D)\to\Omega],
\]
其中 \(0\) 表示 \(E_d^\vee\) 的零截面（它同构于 \(\Bun^{-d}_{\Gm}\)）。

\begin{exercise}\label{ex:coeff-four}
函子 \(\coeff_{d+(2g-2)}\) 可以等同于复合
\[
\Shv(\Bun_G)\xrightarrow{\ p^!_{-d}\ }\Shv(\Bun^{-d}_B)\xrightarrow{\ \Four\ }\Shv(E_d^\vee)
\xrightarrow{\ \text{限制}\ }\Shv(\Sym^{d+(2g-2)}(X)).
\]
\end{exercise}

\begin{fact}\label{fact:connected}
若 \(f\colon Y\to Z\) 是纤维连通的平滑映射，则反常层上的拉回函子是全忠实的。
\end{fact}

\begin{fact}\label{fact:connected-irr}
若 \(f\colon Y\to Z\) 是纤维连通或空的平滑映射，且 \(Y\) 非空，则拉回函子在支撑稠密的不可约反常层上是全忠实的。
\end{fact}

\begin{exercise}\label{ex:fibers}
对任意开子叠 \(U_n\subset\Bun_{\mathrm{PGL}_2}\)，存在只依赖于 \(n\) 与 \(g\) 的整数 \(N\)，
使得对任意 \(d>N\)，映射 \(p_{-d}\colon\Bun^{-d}_B\to\Bun_G\) 在 \(U_n\) 上的纤维连通或空。
\end{exercise}

\begin{proof}[命题 \ref{prop:determined} 的证明]
回忆 \(F\) 由它在 \(U_n\subset\Bun_{\mathrm{PGL}_2}\) 上的限制决定，其中 \(n\) 是只依赖于 \(g\) 的整数
（推论 \ref{cor:extend}）。于是由上述事实与练习，存在只依赖于 \(g\) 的整数 \(N\)，使得对任意 \(d>N\)，
函子 \(p^!_{-d}\) 不会丢失 \(F|_{\Bun^{d+2\Z/2\Z}_{\mathrm{PGL}_2}}\) 的信息。因此，为证明命题 \ref{prop:determined}，
只需从 \(\Four(p^!_{-d}F)\) 沿开浸入
\[
E_d^\vee\setminus 0\subset E_d^\vee
\]
的限制恢复 \(\Four(p^!_{-d}F)\)。

为此注意：\(\Four(p^!_{-d}F)\) 在零截面上的 \(!\)-限制可以等同于 \(\CT^{-d}_*(F)\)；
由于 \(F\) 被假定为尖点，后者为零。
\yizhu{命题 \ref{prop:determined} 的证明中，「零截面上的 \(!\)-限制为零」是关键：它说明
\(\Four(p^!_{-d}F)\) 在 \(E_d^\vee\) 的零截面附近没有「额外」的部分，因此它由开子集
\(E_d^\vee\setminus 0\) 上的限制唯一决定。这一步对应模形式中「常数项为零」的尖点条件。}{原文省略·I：本工程 §2.1 与 §3.2 的衔接}
\end{proof}

\subsection{几何 Langlands 对应}
\srcpage{10}

我们已经看到：若 \(\cF_\sigma\) 存在，则它由 \(\sigma\) 唯一确定。下述结果称为 \(\mathrm{PGL}_2\) 的几何 Langlands 对应。

\begin{theorem}[Drinfeld, Laumon, Frenkel--Gaitsgory--Vilonen]\label{thm:gl2}
对任意不可约 \(\mathrm{SL}_2\)-局部系统 \(\sigma\)，存在 \(\cF_\sigma\in\Shv(\Bun_{\mathrm{PGL}_2})\)，
它反常，且在 \(\Bun_{\mathrm{PGL}_2}\) 的每个连通分支上不可约，并满足
\[
\coeff_d(\cF_\sigma)\simeq\sigma^{(d)}.
\]
此外，\(\cF_\sigma\) 是 Hecke 特征层。
\end{theorem}

为解释「Hecke 特征层」这个概念，考虑模叠 \(\mathrm{Hecke}\)，它分类 \((x\in X,\cE\subset\cE'\subset\cE(x))\)，
其中（模去张量线丛）\(\cE\) 与 \(\cE'\) 是 \(X\) 上的秩 \(2\) 丛，且包含 \(\cE\subset\cE'\subset\cE(x)\) 是严格的。
有对应
\[
\begin{tikzcd}
& \mathrm{Hecke} \arrow[dl,"\overleftarrow{h}"'] \arrow[dr,"\overrightarrow{h}"] & \\
\Bun_G & & \Bun_G\times X
\end{tikzcd}
\]
其中 \(\overleftarrow{h}\) 把 \((x\in X,\cE\subset\cE'\subset\cE(x))\) 送到 \(\cE\)，
\(\overrightarrow{h}\) 把它送到 \((\cE',x)\)。考虑函子
\yizhu{Hecke 叠的严格包含 \(\cE\subset\cE'\subset\cE(x)\) 可以这样理解：\(\cE'/\cE\) 是 \(x\) 处的
一维「跳跃」，而 \(\cE(x)=\cE\otimes\cO(x)\) 是给 \(\cE\) 在 \(x\) 处加一个极点。模去张量线丛后，
这个三元组只依赖 \((\cE,\cE')\) 的相对位置。Hecke 函子 \(H_{\mathrm{std}}\) 的「先 \(!\) 拉回、再 \(*\) 推前」
是几何 Langlands 中 Hecke 算子的标准形式。}{背景直觉·A：Frenkel--Gaitsgory--Vilonen, \emph{On the geometric Langlands conjecture}, J. Amer. Math. Soc. 15 (2002)；B：nLab ``geometric Langlands correspondence''}
\begin{equation}\label{eq:3.2}
H_{\mathrm{std}}:=\overrightarrow{h}_*\,\overleftarrow{h}^!.
\end{equation}
Hecke 特征性质说的是
\begin{equation}\label{eq:3.3}
H_{\mathrm{std}}(\cF_\sigma)\simeq\cF_\sigma\boxtimes\sigma.
\end{equation}

% ============================================================
%  第四讲
% ============================================================
\section{第四讲}

\subsection{Hecke 函子}
\srcpage{11}

现在对一般约化群 \(G\)（在 \(k\) 上分裂且连通）定义 Hecke 性质。设 \(\Gd\) 是 \(\e\) 上的对偶约化群。

我们首先有 \(G\)-丛的概念。

\begin{example}\label{ex:g-bundles}
\(G\)-丛具体表现为：

\begin{itemize}
\item \(G=\mathrm{GL}_n\)：秩 \(n\) 向量丛 \(\cE\)；
\item \(G=\mathrm{SL}_n\)：秩 \(n\) 向量丛 \(\cE\)，并带有 \(\det\cE\simeq\cO\)；
\item \(G=\mathrm{PGL}_n\)：秩 \(n\) 向量丛 \(\cE\)，模去张量线丛；
\item \(G=\mathrm{O}_n\)：秩 \(n\) 向量丛 \(\cE\)，并带有对称非退化形式 \(\cE\otimes\cE\to\cO\)；
\item \(G=\mathrm{Sp}_n\)：秩 \(n\) 向量丛 \(\cE\)，并带有交错非退化形式 \(\cE\otimes\cE\to\cO\)；
\item \(\cdots\)
\end{itemize}
\end{example}

对任意概形 \(S\) 上的 \(G\)-丛 \(P_G\)，有右 \(t\)-正合的对称幺半函子
\yizhu{当 \(S\) 是导出概形时，下面的函子并不 \(t\)-正合；原文的「右 \(t\)-正合」只在
普通概形上成立。译文保留此限定。}{原文省略·A：原文脚注 4}
\[
\Rep(G)\to\QCoh(S),\qquad V\mapsto V_{P_G}.
\]
若 \(V\) 有限维，则 \(V_{P_G}\) 是有限秩向量丛。

我们还有 \(\Gd\)-局部系统的概念，它可以定义为右 \(t\)-正合的对称幺半函子
\yizhu{原文此处写 \(\Rep(G)\to\Shv(X)\)，但按上下文（\(\Gd\)-局部系统、\(V_\sigma\)）应为 \(\Rep(\Gd)\)。
这是原文的笔误，译文照录原文并加此注。}{原文笔误·A：与上一段 \(\Rep(G)\to\QCoh(S)\) 对照}
\[
\Rep(G)\to\Shv(X),\qquad V\mapsto V_\sigma.
\]
若 \(V\in\Rep(\Gd)\) 有限维，则 \(V_\sigma\) 是 lisse 的。

现在对任意 \(V\in\Rep(\Gd)\)，可以定义函子
\[
H_V\colon\Shv(\Bun_G)\to\Shv(\Bun_G\times X),
\]
它推广了上一讲定义的 Hecke 函子（即 \(\Gd=\mathrm{SL}_2\)、\(V\) 为标准表示的情形）。
定义所需的部件是下述等价。

\begin{theorem}[几何 Satake]\label{thm:satake}
对任意 \(x\in X\)，有典范幺半函子
\[
\Rep(\Gd)\to\Shv(L^+_xG\backslash L_xG/L^+_xG),\qquad V\mapsto S_V,
\]
其中 \(L_xG:=G(k((t_x)))\)，\(L^+_xG:=G(k[[t_x]])\)，\(t_x\) 是 \(x\) 附近的局部坐标。
\end{theorem}

\begin{remark}\label{rem:satake-indep}
可以验证 \(S_V\) 不依赖于局部坐标 \(t_x\) 的选取。
\end{remark}

双商
\[
\mathrm{Hecke}^{\mathrm{loc}}_x:=L^+_xG\backslash L_xG/L^+_xG
\]
分类 \(D_x:=\operatorname{Spec} k[[t_x]]\) 上的两个 \(G\)-丛 \(P_G,P'_G\)，以及它们在
\(\overset{\circ}{D}_x\) 上的限制之间的一个同构。我们还可以考虑这个叠的整体版本
\(\mathrm{Hecke}_x\)，它分类 \(X\) 上的两个 \(G\)-丛，以及它们在
\(\overset{\circ}{X}:=X\setminus x\) 上的限制之间的一个同构。对应为
\[
\begin{tikzcd}
& \mathrm{Hecke}_x \arrow[dl,"\overleftarrow{h}"'] \arrow[dr,"\overrightarrow{h}"] \arrow[d,"\pi"] & \\
\Bun_G & \mathrm{Hecke}^{\mathrm{loc}}_x & \Bun_G
\end{tikzcd}
\]

\begin{definition}\label{def:hecke-functor}
定义函子
\[
H_{V,x}\colon\Shv(\Bun_G)\to\Shv(\Bun_G),\qquad
F\mapsto\overrightarrow{h}_*\bigl(\overleftarrow{h}^!(F)\overset{!}{\otimes}\pi^!(S_V)\bigr).
\]
让 \(x\in X\) 变动，类似地定义函子
\[
H_V\colon\Shv(\Bun_G)\to\Shv(\Bun_G\times X),
\]
使得 \(H_V(F)|_{\Bun_G\times\{x\}}\simeq H_{V,x}(F)\)。我们称它们为 \emph{Hecke 函子}。
\end{definition}

\begin{exercise}\label{ex:hecke-pgl2}
取 \(G=\mathrm{PGL}_2\)、\(\Gd=\mathrm{SL}_2\)，并考虑标准表示 \(\mathrm{std}\in\Rep(\Gd)\)。
验证上述 Hecke 函子 \(H_{\mathrm{std}}\) 等价于 \eqref{eq:3.2} 中的函子。
\end{exercise}

Hecke 函子以下述意义交换。对 \(V_1,V_2\in\Rep(\Gd)\)，令 \(H_{V_1,V_2}\) 为复合
\[
\Shv(\Bun_G)\xrightarrow{\ H_{V_2}\ }\Shv(\Bun_G\times X)
\xrightarrow{\ H_{V_1}\ }\Shv(\Bun_G\times X\times X).
\]
则有典范识别
\[
H_{V_1,V_2}\simeq\mathrm{swap}\circ H_{V_2,V_1},
\]
其中 \(\mathrm{swap}\) 是由交换两个 \(X\) 因子诱导的 \(\Shv(\Bun_G\times X\times X)\) 上的对合。
更一般地，给定有限集 \(I\) 与以 \(I\) 为指标集的 \(\Gd\)-表示族 \(V=\{V_i\}_{i\in I}\)，
有良定义函子
\[
H_V\colon\Shv(\Bun_G)\to\Shv(\Bun_G\times X^I),
\]
它由任意顺序复合 \(H_{V_i}\) 得到。

\begin{exercise}\label{ex:hecke-compose}
验证
\[
H_{V_1,x}\circ H_{V_2,x}\simeq H_{V_1\otimes V_2,x}.
\]
\end{exercise}

\begin{question}\label{ques:action}
什么作用在 \(\Shv(\Bun_G)\) 上？
\end{question}

对每个 \(x\in X\)，有 \(\Shv(\Bun_G)\) 上的自函子 \(H_{V,x}\)，但 \(H_V\) 并不是自函子。
有一种形式化的办法，可以从这些 Hecke 函子产生自函子。对有限集 \(I\)、
\(V=\boxtimes_{i\in I}V_i\in\Rep(\Gd)^{\otimes I}\) 以及 \(G\in\Shv(X^I)\)，考虑复合
\begin{equation}\label{eq:4.1}
\Shv(\Bun_G)\xrightarrow{\ H_V\ }\Shv(\Bun_G\times X^I)
\xrightarrow{\ -\overset{!}{\otimes}\mathrm{pr}_2^!(G)\ }\Shv(\Bun_G\times X^I)
\xrightarrow{\ \mathrm{pr}_{1,*}\ }\Shv(\Bun_G).
\end{equation}
这给出函子
\[
\Rep(\Gd)^{\otimes I}\otimes\Shv(X^I)\to\End(\Shv(\Bun_G)),
\]
它把 \(V\boxtimes G\) 送到上述自函子。

\begin{remark}\label{rem:integration}
当 \(I=\{1,2\}\) 时，自函子 \eqref{eq:4.1} 是函子 \(H_{V_1,x_1}\circ H_{V_2,x_2}\)（当 \(x_1\neq x_2\)）
与 \(H_{V_1\otimes V_2,x}\)（当 \(x_1=x_2=x\)）对 \(X^2\) 上的「测度」\(G\) 的「积分」。
\end{remark}

对有限集之间的任意映射 \(\alpha\colon I\to J\)，有函子
\[
\Rep(\Gd)^{\otimes I}\xrightarrow{\ \mathrm{mult}_\alpha\ }\Rep(\Gd)^{\otimes J},
\]
它由 \(\Rep(\Gd)\) 的对称幺半结构给出；还有函子
\[
\Delta^\alpha_*\colon\Shv(X^J)\to\Shv(X^I),
\]
它由映射 \(\Delta_\alpha\colon X^J\to X^I\) 诱导。可以验证下图交换：
\[
\begin{tikzcd}
\Rep(\Gd)^{\otimes I}\otimes\Shv(X^J) \arrow[r,"\mathrm{mult}_\alpha\otimes\mathrm{id}"] \arrow[d,"\mathrm{id}\otimes\Delta^\alpha_*"'] & \Rep(\Gd)^{\otimes J}\otimes\Shv(X^J) \arrow[d] \\
\Rep(\Gd)^{\otimes I}\otimes\Shv(X^I) \arrow[r] & \End(\Shv(\Bun_G)).
\end{tikzcd}
\]
于是有函子
\begin{equation}\label{eq:4.2}
\Rep(\Gd)_{\mathrm{Ran}}:=\operatorname*{colim}_{\alpha\colon I\to J}
\Rep(\Gd)^{\otimes I}\otimes\Shv(X^J)\to\End(\Shv(\Bun_G)),
\end{equation}
其中余极限按映射 \(\alpha\colon I\to J\) 取，并对 \(I\) 反变、对 \(J\) 协变。

注意 \(\Rep(\Gd)_{\mathrm{Ran}}\) 带有（对称）幺半结构：把 \(I\xrightarrow{\alpha}J\) 与
\(I'\xrightarrow{\alpha'}J'\) 的项张量成 \(I\sqcup I'\xrightarrow{(\alpha,\alpha')}J\sqcup J'\) 的项。
可以验证 \eqref{eq:4.2} 与两边的幺半结构相容。
\yizhu{\(\Rep(\Gd)_{\mathrm{Ran}}\) 的「Ran」指 Ran 空间：它把 \(X\) 的所有有限子集放在一起，
允许这些点碰撞。\(\Shv(X^I)\) 的因子记录 Hecke 修改发生的位置；当两个点碰撞时，两个 Hecke 函子复合，
对应 \(\Rep(\Gd)\) 的张量积。这个构造的用途在 \S4.2 显现：它让所有 Fourier 系数有同一个「家」。}{背景直觉·A：Gaitsgory--Raskin, GLC I, arXiv:2405.03599, \S2（\(\Rep(\Gd)_{\mathrm{Ran}}\) 在 \(\Shv(\Bun_G)\) 上的作用）}

总之，把所有 Hecke 函子放在一起，得到幺半函子
\[
\Rep(\Gd)_{\mathrm{Ran}}\to\End(\Shv(\Bun_G)).
\]

\subsection{Hecke 性质}
\srcpage{14}

\begin{question}\label{ques:ran}
\(\Rep(\Gd)_{\mathrm{Ran}}\) 能为我们做什么？
\end{question}

我们给出两个答案：

\begin{itemize}
\item 它把所有 Fourier 系数收进同一个「家」；
\item 它给出特征层的完整而正确的定义。
\end{itemize}

先解释第二点。固定 \(X\) 上的 \(\Gd\)-局部系统 \(\sigma\)，有对称幺半函子
\[
\mathrm{ev}_\sigma\colon\Rep(\Gd)_{\mathrm{Ran}}\to\Vect,
\]
它在 \(I\xrightarrow{=}I\) 这一项上的限制为
\[
\Rep(\Gd)^{\otimes I}\otimes\Shv(X^I)
\xrightarrow{\ \sigma^{\boxtimes I}\otimes\mathrm{id}\ }
\Shv(X^I)\otimes\Shv(X^I)
\xrightarrow{\ -\overset{!}{\otimes}-\ }
\Shv(X^I)\xrightarrow{\ \Gamma\ }\Vect,
\]
其中回忆：\(X^I\) 上的 \(\Gd\)-局部系统 \(\sigma^{\boxtimes I}\) 被看作对称幺半函子
\(\Rep(\Gd)^{\otimes I}\to\Shv(X^I)\)。

\begin{definition}\label{def:eigensheaf}
以 \(\sigma\) 为特征值的 \emph{Hecke 特征层}是一个 \(\Rep(\Gd)_{\mathrm{Ran}}\)-线性函子
\[
\Vect\to\Shv(\Bun_G),\qquad \one\mapsto F,
\]
其中 \(\Rep(\Gd)_{\mathrm{Ran}}\) 经由 \(\mathrm{ev}_\sigma\) 作用在 \(\Vect\) 上。
我们常常滥用记号，把 Hecke 特征层本身就记作 \(F\)。
\end{definition}

\begin{exercise}\label{ex:sigma-unique}
证明：\(\sigma\) 由对称幺半函子 \(\mathrm{ev}_\sigma\) 唯一确定。
\end{exercise}

\begin{exercise}\label{ex:eigen-hecke}
设 \(F\) 是以 \(\sigma\) 为特征值的 Hecke 特征层。证明
\begin{equation}\label{eq:4.3}
H_V(F)\simeq F\boxtimes V_\sigma.
\end{equation}
\end{exercise}

\begin{remark}\label{rem:generalize}
注意 \eqref{eq:4.3} 推广了 \eqref{eq:3.3}。
\end{remark}

\subsection{几何 Langlands 对应}
\srcpage{14}

\begin{conjecture}\label{conj:gl}
若 \(\sigma\) 是不可约 \(\Gd\)-局部系统（即它不来自真抛物子群），则存在唯一的 Hecke 特征层 \(\cF_\sigma\)，
并带有同构 \(\coeff_0(\cF_\sigma)\simeq\one\)。
\end{conjecture}

\begin{theorem}[Arinkin 等]\label{thm:gl-general}
{\footnotesize\noindent 完整作者名单：Arinkin--Beraldo--Campbell--Chen--Faergeman--Gaitsgory--Lin--Raskin--Rozenblyum（原文以 ``R.'' 代 Raskin）。\par}
若 \(\operatorname{char}(k)=0\)，则猜想 \ref{conj:gl} 成立。此外还有：
\begin{itemize}
\item \(\cF_\sigma\) 是反常的；
\item \(\cF_\sigma\) 是尖点的；
\item \(\cF_\sigma\) 是半单的。更精确地说，设 \(S_\sigma\) 是 \(\sigma\) 的自同构群，则有分解
\[
\cF_\sigma\simeq\bigoplus_{\rho\in\mathrm{Irr}(S_\sigma)}\cF_{\sigma,\rho}^{\oplus\dim\rho},
\]
其中 \(\cF_{\sigma,\rho}\) 是单反常层；
\item \(\cF_\sigma\) 的特征圈是 \([\mathrm{Nilp}]\)，其中 \(\mathrm{Nilp}\subset T^*\Bun_G\) 是整体幂零锥；
\yizhu{特征圈为 \([\mathrm{Nilp}]\) 说明 \(\cF_\sigma\) 的奇异支撑是整个幂零锥 \(\mathrm{Nilp}\subset T^*\Bun_G\)。
这是整体几何 Langlands 的深刻结论，与尖点性并不矛盾：尖点性是关于常数项消没的，而奇异支撑描述的是
微局部行为。}{背景直觉·A：Gaitsgory--Raskin, GLC I, arXiv:2405.03599, §6}
\item 当 \(g>1\) 时，\(\cF_\sigma\) 的通用秩为
\[
\prod_i d_i^{(2d_i-1)(g-1)},
\]
其中 \(d_i\) 是 \(G\) 的指数。
\yizhu{原文称 \(d_i\) 为 \(G\) 的指数。按此约定，\(\mathrm{GL}_2\) 的指数为 \(1,2\)，公式给出
\(1^{1}\cdot 2^{3(g-1)}=2^{3g-3}\)，与已知的 \(\mathrm{GL}_2\) Hecke 特征层通用秩一致；
\(\mathrm{PGL}_2\) 的指数为 \(1\)，公式给出 \(1\)，也与其特征层通用秩一致。因此原文公式无需改动。
注意不要把 \(d_i\) 读成不变多项式的次数：\(\mathrm{PGL}_2\) 的不变多项式次数是 \(2\)，代入会得到
不同的数值，那不是原文的约定。}{术语背景·A：原文公式；Gaitsgory--Raskin, GLC I, arXiv:2405.03599, §6}
\end{itemize}
\end{theorem}

\begin{remark}\label{rem:components}
\(\Bun_G\) 的连通分支可以与 \(\mathrm{Irr}(Z(\Gd))\)（\(\Gd\) 的中心的不可约表示之集合）等同。
反常层 \(\cF_{\sigma,\rho}\) 支撑在由 \(\rho|_{Z(\Gd)}\) 标号的连通分支上。
\end{remark}

% ============================================================
%  第五讲
% ============================================================
\section{第五讲}

\subsection{Whittaker 系数}
\srcpage{15}

我们有如下箴言：

\begin{quote}
\(\Rep(\Gd)_{\mathrm{Ran}}\) 让我们把 \(\Bun_G\) 上一个层的所有 Whittaker 系数「粘」在一起。
\end{quote}

为解释这句话，先对一般 \(G\) 定义 Whittaker 系数。考虑
\[
\Bun^\Omega_N:=\Bun_B\times_{\Bun_T}\{\check\rho(\Omega)\},
\]
其中 \(\check\rho(\Omega):=2\check\rho(\Omega^{1/2})\)，\(\Omega^{1/2}\) 是线丛 \(\Omega\) 的一个固定平方根。
对每个正单根 \(\alpha_i\)，映射 \(N\to\Ga\) 诱导出映射
\[
\phi_i\colon\Bun^\Omega_N\to\Bun^\Omega_{\Ga}\to H^1(X,\Omega)\simeq\A^1.
\]
定义
\[
\phi:=\sum_{\alpha_i}\phi_i\colon\Bun^\Omega_N\to\A^1,
\]
并考虑对应
\[
\begin{tikzcd}
& \Bun^\Omega_N \arrow[dl,"p_N"'] \arrow[dr,"\phi"] & \\
\Bun_G & & \A^1
\end{tikzcd}
\]

与 \(\mathrm{PGL}_2\) 情形一样，我们作如下定义。

\begin{definition}\label{def:whittaker0}
设 \(\coeff_0\colon\Shv(\Bun_G)\to\Vect\) 为由下式给出的函子：
\[
\coeff_0(F):=\mathrm{C}^\bullet_c\bigl(\Bun^\Omega_N,\ p_N^!(F)\overset{!}{\otimes}\phi^!(\exp)\bigr).
\]
类似地，对每个 \(D=\sum\check\lambda_i x_i\)（其中 \(\lambda_i\in\check\Lambda^+\)），用
\(\check\rho(\Omega)(-D)\in\Bun_T\) 产生函子
\[
\coeff_D\colon\Shv(\Bun_G)\to\Vect.
\]
我们称它们为 \emph{Whittaker 系数函子}。
\yizhu{Whittaker 系数的经典原型是自守表示上的 Whittaker 函数：固定一个非平凡特征 \(\psi\) 于单幂根，
取矩阵系数的 \(\psi\)-分量。几何情形把「非平凡特征」换成 \(\exp\) 层（de Rham 的 \(e^t\) 或
\(\ell\)-进的 Artin--Schreier 层），把「积分」换成紧支集上同调。}{背景直觉·A：Frenkel--Gaitsgory--Vilonen, \emph{Whittaker patterns in the geometry of moduli spaces of bundles on curves}, Ann. of Math. 153 (2001)；B：nLab ``Whittaker function''}
\end{definition}

\begin{example}\label{ex:gm-whittaker}
对 \(G=\Gm\)，\(D\) 可以是任意除子，函子
\[
\coeff_D\colon\Shv(\Bun_{\Gm})\to\Vect
\]
就是取 \(\cO(-D)\) 处的 \(!\)-纤维。
\end{example}

\begin{example}\label{ex:pgl2-whittaker}
对 \(G=\mathrm{PGL}_2\)，这些函子回到前面讲次中用同样记号表示的函子。
\end{example}

下述结果称为几何 Casselman--Shalika 公式。

\begin{theorem}[Frenkel--Gaitsgory--Vilonen]\label{thm:casselman-shalika}
对 \(D=\sum\check\lambda_i x_i\)（\(\lambda_i\in\check\Lambda^+\)），有
\[
\coeff_D(F)\simeq\coeff_0\bigl(H_{\boxtimes V_{\check\lambda_i},x}(F)\bigr),
\]
其中回忆 \(H_{\boxtimes V_{\check\lambda_i},x}\) 是任意顺序复合 \(H_{V_{\check\lambda_i},x_i}\)。
\end{theorem}

注意有对象
\[
\boxtimes_i\bigl(V_{\check\lambda_i}\otimes\delta_{x_i}\bigr)
\in\Rep(\Gd)^{\otimes I}\otimes\Shv(X^I),
\]
它给出 \(\Rep(\Gd)_{\mathrm{Ran}}\) 中的一个对象。上述定理说的是：\(\coeff_D\) 等价于先作用该对象的 Hecke 作用，
再复合 \(\coeff_0\)。

\begin{definition}\label{def:normalized}
设 \(\cF_\sigma\) 是以 \(\sigma\) 为特征值的 Hecke 特征层。称 \(F\) 是\emph{归一化的}，
若它带有同构
\[
\coeff_0(F)\simeq\one.
\]
\end{definition}

设 \(\cF_\sigma\) 是归一化的 Hecke 特征层。由练习 \ref{ex:eigen-hecke}，有
\[
H_{\boxtimes V_{\check\lambda_i},x}(\cF_\sigma)
\simeq\cF_\sigma\otimes\bigotimes_i(V_{\check\lambda_i,\sigma})_{x_i}.
\]
由上述定理，这个条件蕴含
\[
\coeff_D(\cF_\sigma)
\simeq\coeff_0(\cF_\sigma)\otimes\bigotimes_i(V_{\check\lambda_i,\sigma})_{x_i}
\simeq\bigotimes_i(V_{\check\lambda_i,\sigma})_{x_i},
\]
其中最后一个同构是因为 \(\cF_\sigma\) 归一化。注意这推广了 \eqref{eq:3.1}。

现在考虑下述函子
\begin{equation}\label{eq:5.1}
\Rep(\Gd)_{\mathrm{Ran}}\otimes\Shv(\Bun_G)
\xrightarrow{\ \text{Hecke 作用}\ }\Shv(\Bun_G)
\xrightarrow{\ \coeff_0\ }\Vect.
\end{equation}
上述定理表明，这个函子可以看作把所有 Whittaker 系数装配成一个族。范畴
\(\Rep(\Gd)_{\mathrm{Ran}}\) 是自对偶的，配对函子为
\[
\Rep(\Gd)_{\mathrm{Ran}}\otimes\Rep(\Gd)_{\mathrm{Ran}}
\xrightarrow{\ -\star-\ }\Rep(\Gd)_{\mathrm{Ran}}\xrightarrow{\ \Gamma\ }\Vect,
\]
其中 \(-\star-\) 是 \(\Rep(\Gd)_{\mathrm{Ran}}\) 上的对称幺半结构，\(\Gamma:=\Hom(\mathrm{triv},-)\)。
于是函子 \eqref{eq:5.1} 可以改写成函子
\[
\coeff^{\mathrm{ult}}\colon\Shv(\Bun_G)\to\Rep(\Gd)_{\mathrm{Ran}},
\]
它由下述公式刻画：
\[
\Gamma\bigl(\coeff^{\mathrm{ult}}(F)\star G\bigr)
\simeq\coeff_0(G\cdot F),\qquad F\in\Shv(\Bun_G),\ G\in\Rep(\Gd)_{\mathrm{Ran}}.
\]
这里 \(G\cdot F\) 是 \(G\) 在 \(F\) 上的作用。

\begin{remark}\label{rem:q-expansion}
应当把 \(\coeff^{\mathrm{ult}}\) 看作自守层的「\(q\)-展开」的最佳版本。
\end{remark}

\begin{theorem}[Beraldo, Frenkel--Gaitsgory--Vilonon]\label{thm:coeff-ff}
对 \(G=\mathrm{GL}_n\) 或 \(\mathrm{PGL}_n\)，函子 \(\coeff^{\mathrm{ult}}\) 在
\(\Shv(\Bun_G)_{\mathrm{cusp}}\) 上是全忠实的。
\yizhu{定理 \ref{thm:coeff-ff} 的原文把最后一位作者拼成 ``Vilonon''，应为 ``Vilonen''。
这是排印错误，不影响数学内容。}{原文笔误·A：作者名}
\end{theorem}

对 \(\mathrm{PGL}_2\)，证明模仿 Fourier 反演；对 \(\mathrm{GL}_n\)，证明模仿
Piatetski-Shapiro 与 Shalika 的工作。

\subsection{范畴化的几何 Langlands 等价}
\srcpage{17}

现在我们在 de Rham 情形工作。考虑 \(\mathrm{LS}_{\Gd}=\mathrm{LS}^{\dR}_{\Gd}\)，
即 \(X\) 上（de Rham）\(\Gd\)-局部系统的模叠。存在典范对称幺半函子
\[
\Loc\colon\Rep(\Gd)_{\mathrm{Ran}}\to\QCoh(\mathrm{LS}_{\Gd}),
\]
它把 \(V\otimes\delta_x\in\Rep(\Gd)\otimes\DMod(X)\) 送到 \(\mathrm{ev}_x^*(V)\)。这里
\[
\mathrm{ev}_x\colon\mathrm{LS}_{\Gd}\to\mathrm{LS}_{\Gd}(D_x)\simeq B\Gd
\]
是赋值映射，并且我们把 \(\QCoh(B\Gd)\) 与 \(\Rep(\Gd)\) 等同。

\begin{exercise}\label{ex:ev-loc}
对 \(\sigma\in\mathrm{LS}_{\Gd}\)，把函子
\[
\mathrm{ev}_\sigma\colon\Rep(\Gd)_{\mathrm{Ran}}\to\Vect
\]
等同于复合
\[
\Rep(\Gd)_{\mathrm{Ran}}\xrightarrow{\ \Loc\ }\QCoh(\mathrm{LS}_{\Gd})
\xrightarrow{\ (-)^!_\sigma\ }\Vect.
\]
\end{exercise}

\begin{theorem}[Lurie, Gaitsgory--Rozenblyum]\label{thm:loc-adjoint}
函子 \(\Loc\) 有全忠实的右伴随
\[
\QCoh(\mathrm{LS}_{\Gd})\hookrightarrow\Rep(\Gd)_{\mathrm{Ran}}.
\]
\end{theorem}

\begin{theorem}[Drinfeld--Gaitsgory]\label{thm:action-factors}
\(\Rep(\Gd)_{\mathrm{Ran}}\) 在 \(\DMod(\Bun_G)\) 上的作用经 \(\Loc\) 分解。
\end{theorem}

\begin{corollary}\label{cor:coeff-factors}
函子 \(\coeff^{\mathrm{ult}}\) 分解为
\[
\DMod(\Bun_G)\xrightarrow{\ \LG_{\mathrm{coarse}}\ }\QCoh(\mathrm{LS}_{\Gd})
\hookrightarrow\Rep(\Gd)_{\mathrm{Ran}}.
\]
\end{corollary}

范畴化几何 Langlands 的想法是：函子 \(\LG_{\mathrm{coarse}}\) 「几乎」是等价。注意它不可能是等价，
因为对非交换的 \(G\)，它把常层送到 \(0\)。

\begin{theorem}[范畴化几何 Langlands，第一版]\label{thm:categorical-1}
函子 \(\LG_{\mathrm{coarse}}\) 诱导等价
\[
\DMod(\Bun_G)_{\mathrm{cusp}}\simeq\QCoh(\mathrm{LS}^{\mathrm{irred}}_{\Gd}),
\]
其中 \(\mathrm{LS}^{\mathrm{irred}}_{\Gd}\subset\mathrm{LS}_{\Gd}\) 是不可约 \(\Gd\)-局部系统的轨迹。
\end{theorem}

\begin{corollary}\label{cor:coeff-ff-cusp}
\(\coeff^{\mathrm{ult}}\) 在尖点子范畴上是全忠实的。
\end{corollary}

Arinkin--Gaitsgory 对 \(\LG_{\mathrm{coarse}}\) 还有一个修正。回忆：对足够好的叠 \(Y\)，有
\[
\QCoh(Y)\simeq\Ind(\Perf(Y)),
\]
其中 \(\Perf(Y)\) 是局部有界、有限秩投射模复形构成的全子范畴。我们还可以考虑
\[
\IndCoh(Y):=\Ind(\Coh(Y)),
\]
其中 \(\Coh(Y)\) 是局部有界、有限生成模复形构成的范畴。有包含
\[
\Perf(Y)\subset\Coh(Y),
\]
当 \(Y\) 有奇点时，这个包含是严格的。

现在可以定义子范畴
\[
\Coh_{\mathrm{Nilp}}(\mathrm{LS}_{\Gd})\subset\Coh(\mathrm{LS}_{\Gd}),
\]
它由所有抛物子群 \(\check P\) 及其 Levi 商 \(\check M\) 的下述像生成：
\[
\Perf(\mathrm{LS}_{\check M})\xrightarrow{\ q^* \ }\Perf(\mathrm{LS}_{\check P})
\xrightarrow{\ p_* \ }\Coh(\mathrm{LS}_{\Gd}).
\]
这里 \(q\colon\mathrm{LS}_{\check P}\to\mathrm{LS}_{\check M}\) 是有限 Tor 幅度的映射，
\(p\colon\mathrm{LS}_{\check P}\to\mathrm{LS}_{\Gd}\) 是固有映射，因此上述函子有定义。
\yizhu{\(\IndCoh_{\mathrm{Nilp}}\) 的下标 ``Nilp'' 指幂零奇异支撑。Arinkin--Gaitsgory 的修正动机是：
\(\Bun_G\) 一般有奇点，\(\QCoh\) 不是正确的目标；\(\IndCoh\) 加上幂零奇异支撑条件才是。
Geometric Langlands 的最终陈述因此写成 \(\DMod(\Bun_G)\simeq\IndCoh_{\mathrm{Nilp}}(\mathrm{LS}_{\Gd})\)，
而不是与 \(\QCoh\) 的等价。}{背景直觉·A：Arinkin--Gaitsgory, \emph{Singular support of coherent sheaves}; B：nLab ``geometric Langlands correspondence''}

\begin{remark}\label{rem:cohnilp}
子范畴 \(\Coh_{\mathrm{Nilp}}(\mathrm{LS}_{\Gd})\subset\Coh(\mathrm{LS}_{\Gd})\)
可以用凝聚层的奇异支撑理论内蕴地定义。
\end{remark}

\begin{theorem}\label{thm:LG}
存在唯一函子
\[
\LG\colon\DMod(\Bun_G)\to\IndCoh_{\mathrm{Nilp}}(\mathrm{LS}_{\Gd}),
\]
它满足一个技术条件\yizhu{定理 \ref{thm:LG} 的技术条件是：把源范畴中的紧对象送到目标 \(t\)-结构中有下界的对象。
这是原文脚注 6 的内容；译文中置于脚注，正文不重复。}{记号·A：原文脚注 6}，并使得下图交换：
\[
\begin{tikzcd}
\DMod(\Bun_G) \arrow[r,"\LG"] \arrow[dr,"\LG_{\mathrm{coarse}}"'] & \IndCoh_{\mathrm{Nilp}}(\mathrm{LS}_{\Gd}) \arrow[d,"\Psi"] \\
& \QCoh(\mathrm{LS}_{\Gd}),
\end{tikzcd}
\]
其中 \(\Psi\) 是包含 \(\Coh_{\mathrm{Nilp}}\subset\Coh\subset\QCoh\) 的 ind-延拓。
\end{theorem}

\begin{theorem}[范畴化几何 Langlands，最终版]\label{thm:categorical-ultimate}
函子 \(\LG\) 是等价。
\end{theorem}

\begin{remark}\label{rem:proof}
证明用到 de Rham 情形的特殊性，例如用临界水平的 Kac--Moody 局部化来证明关于
\(\DMod(\Bun_G)\) 的「几何」陈述。更严肃地说，我们用到 de Rham 模叠
\(\mathrm{LS}^{\dR}_{\Gd}\) 的全局（代数）函数很少这一事实。
尽管如此，我们仍可以通过某种意义上的「Riemann--Hilbert」得到 Betti 版本
（以及 \(\operatorname{char}(k)=0\) 时的 \(\ell\)-进版本）。
\end{remark}

\begin{remark}\label{rem:concrete}
上一讲的「具体」定理 \ref{thm:gl-general} 实际上用到了全部范畴断言。
\end{remark}

% ============================================================
%  参考文献、译名对照、原文勘误
% ============================================================
\raggedright

\section*{附录 A：参考文献}
\addcontentsline{toc}{section}{附录 A：参考文献}

\subsection*{A 级：原始文献、经典教材、作者本人文献}
\begin{enumerate}
\item S.~Raskin (lecturer), L.~Chen (notes), \emph{Notes for ``Geometric Langlands for Projective Curves''},
July 07--11, 2025, 18 pp. 公开副本：
\url{https://www.math.columbia.edu/~okounkov/Sam_lecture_notes.pdf}.
\item V.~Drinfeld, D.~Gaitsgory, \emph{Geometric constant term functor(s)},
Selecta Math. (N.S.) \textbf{22} (2016), 1881--1951. Zbl 1360.14060.
\item B.~C.~Ngo, P.~Polo, \emph{Résolutions de Demazure affines et formule de Casselman--Shalika géométrique},
arXiv:math/0005022（预印本，承诺发表于 J. Algebraic Geom.；底本 §5.1 将公式归给 Frenkel--Gaitsgory--Vilonen）.
\item E.~Frenkel, D.~Gaitsgory, K.~Vilonen, \emph{Whittaker patterns in the geometry of moduli spaces of bundles on curves},
Ann. of Math. (2) \textbf{153} (2001), 699--748.
\item E.~Frenkel, D.~Gaitsgory, K.~Vilonen, \emph{On the geometric Langlands conjecture},
J. Amer. Math. Soc. \textbf{15} (2002), 367--417; arXiv:math/0012255.
\item D.~Gaitsgory, S.~Raskin, \emph{Proof of the geometric Langlands conjecture I: construction of the functor},
arXiv:2405.03599.
\item D.~Arinkin, D.~Gaitsgory, \emph{Singular support of coherent sheaves and the geometric Langlands conjecture},
Selecta Math. (N.S.) \textbf{21} (2015), 1--199.
\item P.~Deligne, \emph{Équations différentielles à points singuliers réguliers},
Lecture Notes in Math. \textbf{163}, Springer, 1970（Riemann--Hilbert 对应与正则奇点）.
\end{enumerate}

\subsection*{B 级：可核对的百科、正式讲义}
\begin{enumerate}
\item T.~Feng, \emph{Geometric class field theory}, 讲义,
\url{https://math.berkeley.edu/~fengt/2GeometricCFT.pdf}.
\item nLab, ``local system'', Last revised 2025-04-10,
\url{https://ncatlab.org/nlab/show/local+system}.
\item nLab, ``perverse sheaf'', \url{https://ncatlab.org/nlab/show/perverse+sheaf}.
\item nLab, ``geometric Langlands correspondence'',
\url{https://ncatlab.org/nlab/show/geometric+Langlands+correspondence}.
\item nLab, ``Ran space'', \url{https://ncatlab.org/nlab/show/Ran+space}.
\item nLab, ``Whittaker function'', \url{https://ncatlab.org/nlab/show/Whittaker+function}.
\end{enumerate}

\subsection*{I 级：本库既有译注、术语表、资料卡}
\begin{enumerate}
\item 本库既有卡 \texttt{glc.concept.deligne-gl1-existence}（Deligne GL\(_1\) 定理，本质唯一）.
\item 本库既有卡 \texttt{glc.concept.unramified-global-gl}（非分歧整体几何 Langlands 猜想）.
\item 本库既有卡 \texttt{glc.concept.cuspidal-de-rham}（尖点 de Rham 等价）.
\item 本库既有卡 \texttt{gaits04.concept.eigensheaf}（Hecke 特征层）.
\item 本库既有卡 \texttt{bdh.concept.hecke-eigensheaf}（D-模侧的 Hecke 特征层）.
\end{enumerate}

\section*{附录 B：译名对照}
\addcontentsline{toc}{section}{附录 B：译名对照}

\begin{center}
\small
\begin{tabular}{lll}
\hline
英文 & 中文 & 备注 \\
\hline
local system & 局部系统 & 与 lisse 层同义 \\
lisse & lisse & 保留原文，不译 \\
perverse sheaf & 反常层 & 严禁「错致层」 \\
cuspidal & 尖点 & 作形容词时写「尖点层」 \\
constant term & 常数项 & \\
Whittaker coefficient & Whittaker 系数 & 人名不译 \\
Hecke eigensheaf & Hecke 特征层 & \\
geometric Satake equivalence & 几何 Satake 等价 & \\
affine Grassmannian & 仿射 Grassmann 簇 & \\
loop group & 圈群 & \\
positive loop group & 正圈群 & \\
local Hecke stack & 局部 Hecke 叠 & \\
ind-coherent sheaves & 归纳凝聚层 & IndCoh 保留记号 \\
perfect complex & 完美复形 & \\
quasicoherent & 拟凝聚 & \\
six-functor formalism & 六函子形式化 & \\
correspondence & 对应 & \\
constructible sheaf & 可构造层 & \\
flat connection & 平坦联络 & \\
D-module & D-模 & \\
nilpotent singular support & 幂零奇异支撑 & \\
singular support & 奇异支撑 & \\
microlocal & 微局部 & \\
microstalk & 微茎 & \\
cotangent bundle & 余切丛 & \\
Borel subgroup & Borel 子群 & 人名不译 \\
maximal torus & 极大环面 & \\
reductive group & 约化群 & \\
unipotent radical & 单幂根 & \\
parabolic subgroup & 抛物子群 & \\
Levi subgroup & Levi 子群 & \\
Langlands dual group & Langlands 对偶群 & \\
Eisenstein series & Eisenstein 级数 & \\
Kac--Moody localization & Kac--Moody 局部化 & \\
critical level & 临界水平 & \\
spherical vector & 球向量 & \\
automorphic form & 自守形式 & \\
automorphic representation & 自守表示 & \\
function field & 函数域 & \\
\hline
\end{tabular}
\end{center}

\section*{附录 C：原文勘误表}
\addcontentsline{toc}{section}{附录 C：原文勘误表}

下表只列本译本核对过的、影响理解的原文笔误；拼写类错误集中在末行，不逐条列入。

\begin{center}
\small
\begin{tabular}{>{\raggedright\arraybackslash}p{1.2cm}>{\raggedright\arraybackslash}p{4.4cm}>{\raggedright\arraybackslash}p{3.8cm}>{\raggedright\arraybackslash}p{3.8cm}}
\hline
页 & 原文 & 应为 & 依据 \\
\hline
4 & \(AJ_d\) 的纤维是 \(\Gamma(X,L)\setminus 0\) & \(\mathbb P\Gamma(X,L)\) & 按除子定义；Feng 讲义 \\
4 & preverse & perverse & 拼写 \\
6 & \(B H^1(X,L)\)；\(H^1(X,L)\) 可缩 & \(R\Gamma(X,L)[1]\)；\(d>2g-2\) 时为 \(B H^0(X,L)\) & 脚注与 \(H^1=0,H^0\neq0\) \\
6 & \(0\to L\to B\to \cO\to 0\) & \(0\to L\to \cE\to \cO\to 0\) & 正合列中 \(B\) 为 \(\cE\) 之误 \\
7 & holonomic modular form & holomorphic modular form & 上下文 \\
9 & \(\sigma^{(d)}\) perverse but not lisse & \(\sigma^{(d)}\) lisse；\(\cF_\sigma\) perverse but not lisse & 定义 \\
12 & \(\Rep(G)\to\Shv(X)\) & \(\Rep(\Gd)\to\Shv(X)\) & \(\Gd\)-局部系统的定义 \\
16 & Vilonon & Vilonen & 作者名 \\
多处 & irrducible, uniquelly, Geothendieck, numbder, restiction, Frobenious, archemedean, mearsure, swaping, wher, cmoplexes, prametrizing & 对应正字 & 拼写 \\
\hline
\end{tabular}
\end{center}

\section*{附录 D：底本与版权说明}
\addcontentsline{toc}{section}{附录 D：底本与版权说明}

底本为 Sam Raskin 讲授、Lin Chen 笔记的公开讲义，原 PDF 共 18 页，PDF 页与印刷页一致。
本译本是中文学术译注，不替代原文；原文著作权归作者所有。译本中的译注、术语说明与勘误表
由译者撰写，按本库规范以「改写」（rewritten）方式登记。引用本译本时请同时引用底本。


\end{document}
